\section{A uniform cubic estimate for the reduced action}
\label{sec:v79-cubic}
The second derivatives of the internal contact map need not be
uniformly bounded near grazing. They are not needed to estimate the
third derivative of the reduced roof. Stationarity cancels them.
Throughout this section derivatives are evaluated at actual regular
physical trajectories; no continuation through a singularity is made.

\begin{lemma}[Cubic differentiation of a stationary action]
\label{lem:v79-stationary-cubic}
Let $\mathcal S(q,z)$ be smooth and suppose that
$\partial_z\mathcal S(q,z(q))=0$, with invertible internal Hessian.
Set $F(q)=\mathcal S(q,z(q))$ and
$U_u=(u,Dz(q)u)$. Then
\begin{equation}\label{eq:v79-envelope-cubic}
 D^3F(q)[u,v,w]=D^3\mathcal S(q,z(q))[U_u,U_v,U_w].
\end{equation}
The equality is pointwise and requires no bound for the inverse
internal Hessian or for $D^2z$.
\end{lemma}
\begin{proof}
The differentiated stationary equation says
$D^2\mathcal S[U_u,(0,\zeta)]=0$ for every internal vector $\zeta$.
Consequently $D^2F[u,v]=D^2\mathcal S[U_u,U_v]$.
Differentiate this identity in direction $w$. The two terms in which
$U_u$ or $U_v$ is differentiated are contractions of $D^2\mathcal S$
with an internal vector and vanish by the same stationary equation.
The remaining term is \eqref{eq:v79-envelope-cubic}. This calculation
is local, where the implicit function theorem gives a smooth $z(q)$.
\end{proof}

Use arclength at every contact, including $s_0=R\alpha_0$ and
$s_m=R\alpha_m$, and write
\[
 \mathcal S(s_0,\ldots,s_m)
       =\sum_{j=0}^{m-1}|Q_{j+1}(s_{j+1})-Q_j(s_j)|.
\]
The endpoints are linear functions of $q=(\alpha_0,\alpha_m)$.
At a reflection the internal arclength derivative of $\mathcal S$
is zero. Its internal Hessian is $D_cAD_c$, with $A$ as in
\eqref{eq:v79-optical-matrix}; this is positive definite at every
regular word. In particular Lemma~\ref{lem:v79-stationary-cubic}
applies even when some positive incidence is arbitrarily small.

\begin{theorem}[Collision-uniform third derivatives]
\label{thm:v79-uniform-cubic}
Let $C_Q$ be the explicit constant in
\eqref{eq:v78-contact-derivative}, and put
\begin{equation}\label{eq:v79-cubic-constant}
 C_3=\frac6{\tau_-^2}+\frac3{R_-\tau_-}+\frac1{R_-^2},
 \qquad T_* =\frac{16 C_3 C_Q^3}{1-q_0^3}.
\end{equation}
For every regular physical word, its reduced roof in endpoint angles
satisfies
\begin{equation}\label{eq:v79-uniform-cubic}
 |D^3F(q)[u,v,w]|\le T_*|u|\,|v|\,|w|.
\end{equation}
The constant is independent of $m,R,n,k$, all positive internal
incidences and all positive third-disk clearances. In particular,
on any regular endpoint curve $\gamma$,
\begin{equation}\label{eq:v79-hessian-motion}
 \|K(\gamma(b))-K(\gamma(a))\|
                       \le T_*\operatorname{length}(\gamma).
\end{equation}
This does not assert a common continuation neighborhood for the words.
\end{theorem}
\begin{proof}
For $N(z)=|z|$ at a nonzero vector of length at least $\tau_-$,
\[
 \|DN\|\le1,\qquad \|D^2N\|\le\tau_-^{-1},\qquad
                           \|D^3N\|\le6\tau_-^{-2}.
\]
A circular arclength parametrization has derivative norms
$1,R^{-1},R^{-2}$ through order three. The chain rule for one flight
therefore bounds its third derivative on contact vectors $U,V,W$ by
\[
 C_3(|U_j|+|U_{j+1}|)(|V_j|+|V_{j+1}|)(|W_j|+|W_{j+1}|).
\]
These are derivatives of the action on the product of circles, not
third derivatives of the contact map.

For a lifted endpoint direction $u$, Lemma~\ref{lem:v78-inverse}
gives $|(U_u)_j|\le C_Qq_0^{d_j}|u|$. Indeed a circle has unit speed
in arclength, so its contact-position derivative has precisely the
absolute value of its arclength derivative. The two contact depths
of flight $j$ are at least $d'_j$. Thus this flight contributes at
most $8C_3C_Q^3q_0^{3d'_j}|u||v||w|$ to
\eqref{eq:v79-envelope-cubic}. There are at most two flights at every
depth. Summation of the geometric series proves
\eqref{eq:v79-uniform-cubic}, including $m=1$, when there are no
internal variables. Integrating $DK=D^3F$ along $\gamma$ proves
\eqref{eq:v79-hessian-motion}.
\end{proof}

\begin{corollary}[Uniform angular Jacobian variation]
\label{cor:v79-angular-jacobian}
Let $\Phi_h$ be the canonical forward flow of
Theorem~\ref{thm:v77-field}. On its surviving exact-label trajectories,
with $C_X=\sqrt{2\Lambda_\circ}/\lambda_\circ$,
\begin{equation}\label{eq:v79-chi-motion}
 |\chi(\Phi_hx)-\chi(x)|\le T_*C_X\sqrt h,\qquad
 \left|\log\frac{\chi(\Phi_hx)}{\chi(x)}\right|
                         \le\frac{T_*C_X}{\lambda_\circ}\sqrt h.
\end{equation}
In endpoint coordinates the Jacobian of the flow relative to $dq$ is
\begin{equation}\label{eq:v79-q-jacobian}
               J_h^q(x)=\frac{\chi(x)}{\chi(\Phi_hx)}.
\end{equation}
The physical source Jacobian is different and is computed below.
\end{corollary}
\begin{proof}
In dimension two $\chi=(e^\perp)^{\mathsf T}Ke^\perp$, and $e$ is
constant along the canonical flow. Apply
\eqref{eq:v79-hessian-motion}, the forward length bound, and
$\chi\ge\lambda_\circ$. The map $(F,\vartheta)$ has Jacobian $\chi$
and conjugates the flow to $(t,\theta)\mapsto(t+h,\theta)$.
Ordinary change of variables proves \eqref{eq:v79-q-jacobian}.
No estimate on a derivative of the source density was used.
\end{proof}
